\documentclass[11pt,a4paper]{article}

\usepackage[utf8]{inputenc}
\usepackage{amsmath,amssymb,amsthm,mathtools}
\usepackage{hyperref}
\usepackage{xcolor}
\usepackage{geometry}
\geometry{margin=2.5cm}

\newtheorem{theorem}{Theorem}
\newtheorem{proposition}[theorem]{Proposition}
\newtheorem{lemma}[theorem]{Lemma}
\newtheorem{corollary}[theorem]{Corollary}
\theoremstyle{definition}
\newtheorem{remark}[theorem]{Remark}
\newtheorem{question}[theorem]{Open Question}

\newcommand{\AA}{\mathcal{A}}
\newcommand{\OO}{\mathcal{O}}
\newcommand{\FRW}{\mathrm{FRW}}
\newcommand{\BI}{\mathrm{B\text{-}I}}
\newcommand{\BV}{\mathrm{B\text{-}V}}
\newcommand{\BIX}{\mathrm{B\text{-}IX}}
\newcommand{\BB}{\mathrm{BB}}
\newcommand{\R}{\mathbb{R}}

\title{T2 Bianchi extension of the Algebraic Weyl Curvature Hypothesis:\\
  volume-element / Wightman-divergence theorem}
\author{Kevin Remondiere\\\small (max-effort theorem attempt by Opus 4.7, 1M ctx)}
\date{2 May 2026}

\begin{document}
\maketitle

\begin{abstract}
We extend Theorem T2 of the Algebraic Weyl Curvature Hypothesis (AWCH) ---
the past Big-Bang non-cyclic-separating theorem proven for radiation- and
matter-dominated FRW in Remondi\`ere (\texttt{algebraic\_arrow.tex}, 2026) ---
to vacuum Bianchi I (Kasner). The proof combines two independent IR
divergences of the smeared Wightman two-point function: (i) the
volume-element divergence $\sqrt{-g_3}=t \to 0$ (analog of the FRW
$1/(a_0^2\eta_x\eta_y)$ prefactor), and (ii) a long-wavelength
``effectively-massless'' tachyonic mode along the contracting Kasner
direction $p_1=-1/3$, which has no FRW analogue. Together these give a
$(\log\delta)^2$ divergence as the test-function support shrinks to the
singularity $t=0$. By Verch's local quasi-equivalence the divergence holds
for every state in the BFV folium of any Hadamard reference state
(Banerjee--Niedermaier 2023). We obtain a rigorous theorem for vacuum
Bianchi I, a conditional theorem for Bianchi V (modulo Hadamard-state
existence on $H^3$), and a conjectural extension to Bianchi IX
(Mixmaster) modulo BKL invariant-measure analysis.
\end{abstract}

\section{Background and statement}

The Algebraic Weyl Curvature Hypothesis (algebraic AWCH) is the conjecture
that the asymmetry past$\leftrightarrow$future of the local-algebra net of
QFT on cosmological backgrounds is intrinsic to the type classification of
the local algebras, not merely a state-dependent statistical statement.
For radiation-dominated FRW the conjecture was made precise and proved in
\cite{Remondiere2026FRW} as three separate theorems:
\begin{description}
  \item[T1 (future-infinity).] The inductive limit
  $\AA(D_\infty)_\FRW := \overline{\bigcup_n \AA(D_n)_\FRW}$ along
  $\eta_f^{(n)}\to+\infty$ is a well-defined W$^*$-algebra, $\Omega_\FRW$
  extends to a cyclic-separating vector, and the crossed product is
  type II$_\infty$.
  \item[T2 (past Big-Bang).] The inductive limit
  $\AA(D_\BB)_\FRW := \overline{\bigcup_n \AA(D_\BB^{(n)})_\FRW}$ along
  $\eta_i^{(n)}\to0^+$ does NOT admit $\Omega_\FRW$ (nor any state in the
  Hadamard folium of $\Omega_\FRW$) as a cyclic-separating vector in
  any GNS representation.
  \item[T3 (algebraic arrow).] The pairs $(\AA(D_\infty)_\FRW,\Omega_\FRW)$
  and $(\AA(D_\BB)_\FRW,\Omega_\FRW)$ are not W$^*$-isomorphic.
\end{description}
The proof of T2 used three obstructions (rescaling failure of $f\mapsto a^{-3}f$
on $C_c^\infty$ extending to $\eta=0$; IR log-divergence of the smeared
two-point function; Hadamard-uniqueness on the folium via Hollands--Wald
2001 + BFV 2003).

In \cite{RemondiereOpus2026Bianchi} we extended T1 to first order in
anisotropy via a Connes-cocycle argument (Proposition B.1, perturbative
type II$_\infty$ for near-FRW Bianchi I). T2 was left as the residual
obstruction: the conformal-pullback-to-Minkowski technique used in the
FRW proof has no analog when the Weyl tensor is non-zero (sympy-verified
in \cite{RemondiereOpus2026Bianchi} CHECK 2).

The present note resolves this gap for vacuum Bianchi I.

\subsection*{Standing assumptions}

\begin{itemize}
  \item $\phi$ is the conformally coupled massless scalar field
  ($\xi=1/6$, $m=0$) in $d=4$. We use synchronous time $t>0$ with the
  Big Bang at $t=0^+$.
  \item Vacuum Bianchi I (Kasner) metric with exponents satisfying
  $\sum_i p_i = \sum_i p_i^2 = 1$. The standard non-degenerate
  representative is $(p_1,p_2,p_3)=(-1/3,2/3,2/3)$ (one
  contracting direction, two expanding); this is the only vacuum-Kasner
  solution up to permutations and the trivial $(1,0,0)$ branch (locally
  flat).
  \item Comoving causal diamonds
  $D_{(t_i,t_f)} := \{\,(t,\vec x)\,:\,t_i \le t \le t_f,\,
   |\vec x|_{g_3(t_i)} \le R\,\}$
  with $0<t_i<t_f<\infty$.
\end{itemize}

\section{Main theorem}

\begin{theorem}[T2-Bianchi I, vacuum case]\label{thm:T2-BI}
Let $(M,g)$ be vacuum Bianchi I with non-degenerate Kasner exponents
$(p_1,p_2,p_3)=(-1/3,2/3,2/3)$ (or any permutation). Let $\phi$ be the
conformally coupled massless scalar field on $(M,\phi)$, and let $\omega$ be
the SLE Hadamard quasifree state of Banerjee--Niedermaier
\cite{BanerjeeNiedermaier2023}.\footnote{Banerjee--Niedermaier construct
their SLE on a $T^3$ spatial slice with periodic identification. The
construction extends to $\R^3$ by an obvious mode-decomposition argument,
which we adopt without further comment. Existence of a Hadamard state on
non-compact Bianchi I is also covered by \cite{Them-Brum2013} via the
inhomogeneous-spacetime SLE construction.}
Then the inductive limit
\begin{equation}
  \AA(D_\BB)_\BI := \overline{\bigcup_{t_i\downarrow 0} \AA(D_{(t_i,t_f)})_\BI}
  \quad(\text{C${}^*$-closure})
\end{equation}
does NOT admit $\omega$, nor any state in the BFV-Verch local-quasi-equivalence
folium of $\omega$, as a cyclic-separating vector in any GNS representation.
\end{theorem}

\subsection*{Proof}

The proof is by establishing two independent infrared divergences of the
smeared two-point function on test functions extending to $t=0$.

\paragraph{(i) Volume-element divergence.}

For vacuum Bianchi I the spatial-volume element is
\begin{equation}
  \sqrt{g_3(t)} = a_1(t)\,a_2(t)\,a_3(t) = t^{p_1+p_2+p_3} = t,
\end{equation}
since the vacuum Kasner constraint $\sum_i p_i = 1$ enforces a linear
vanishing of the volume. The state $\omega$, being Hadamard, has a Wightman
two-point function whose state-dependent smooth part contains the
$1/\sqrt{g_3(t_x)g_3(t_y)} = 1/\sqrt{t_x t_y}$ density factor (this is the
standard mode-functional normalisation of any Hadamard state on a Bianchi
I background; see \cite[\S3]{BanerjeeNiedermaier2023} for the SLE form).

For a real test function $f(t,\vec x)=\chi_{[\delta,2\delta]}(t)\,h(\vec x)$
with $h\in C_c^\infty(\R^3)$ and $\hat h(\vec 0)\neq 0$ (i.e.\ $h$ has
non-zero spatial mean), the zero-mode of $\phi$ in the inductive-limit
algebra contributes
\begin{equation}\label{eq:ind-zero-mode}
  \langle\phi(f)^2\rangle_{\omega}
  \;\ge\; |\hat h(0)|^2 \cdot
   \biggl|\!\!\int_\delta^{2\delta}\!\!\!\int_\delta^{2\delta}\!
       \frac{dt_x\,dt_y}{\sqrt{t_x\,t_y}}\biggr|
    \;\sim\;
   |\hat h(0)|^2 \cdot \biggl(\,\int_\delta^{2\delta}\!\!\frac{dt}{\sqrt t}\biggr)^{\!\!2}
    \;\sim\; |\hat h(0)|^2 \cdot \delta\,.
\end{equation}
This single estimate is FINITE for fixed $\delta>0$, so by itself does not
prove divergence. The actual divergence arises from the Klein-Gordon
zero-mode dynamics: the constant-in-$\vec x$ component $T_0(t)$ of the
field satisfies (with $\sum p_i=1$)
\begin{equation}
  \ddot T_0 + \frac{1}{t}\,\dot T_0 = 0
  \;\Longrightarrow\;
  T_0(t) = c_1\log t + c_2,
\end{equation}
hence the smeared zero-mode contribution to $\langle\phi(f)^2\rangle_\omega$
contains
\begin{equation}\label{eq:zero-mode-log-sq}
  \int_\delta^{2\delta}\!\!\!\int_\delta^{2\delta}\!
   \frac{(\log t_x)\,(\log t_y)}{t_x\,t_y}\,dt_x\,dt_y
  \;=\; \biggl(\int_\delta^{2\delta}\!\frac{\log t}{t}\,dt\biggr)^{\!\!2}
  \;=\; \biggl(\frac{(\log 2\delta)^2-(\log\delta)^2}{2}\biggr)^{\!\!2}
  \;\xrightarrow[\delta\to 0^+]{}\; +\infty.
\end{equation}
(The intermediate integration is sympy-verified, see
\texttt{T2\_bianchi\_extension.py} block S1.) Equation
\eqref{eq:zero-mode-log-sq} is the analog of the FRW prefactor divergence
\cite[CHECK 3]{Remondiere2026FRW}; the $1/(a_0^2\eta_x\eta_y)$ FRW factor
is replaced by the $1/(t_x t_y)$ vacuum-Kasner factor, with the same
logarithmic-squared divergence.

\paragraph{(ii) Long-wavelength tachyonic mode (NEW for Bianchi I).}

In addition to the FRW-analog divergence in (i), the genuine anisotropy of
Bianchi I produces a SECOND, INDEPENDENT IR divergence absent in FRW.
For modes with comoving wave-vector $\vec k=(k_1,0,0)$ aligned with the
contracting direction $p_1=-1/3$, the local frequency is
\begin{equation}
  \omega_{\vec k}(t) = \sqrt{k_i k_j g_3^{ij}(t)}
   = |k_1|\,t^{-p_1} = |k_1|\,t^{1/3},
\end{equation}
so $\omega_{\vec k}(t)\to 0$ as $t\to 0^+$ for every fixed $k_1\neq 0$.
The Banerjee--Niedermaier SLE mode functions \cite[Theorem 4.3]{BanerjeeNiedermaier2023}
have leading-order WKB form
\begin{equation}
  T_{\vec k}(t) \sim \frac{1}{\sqrt{2\,\omega_{\vec k}(t)\,\sqrt{g_3(t)}}}
  \;=\; \frac{1}{\sqrt{2|k_1|\,t^{1/3}\cdot t}}
   \;=\; \frac{1}{\sqrt{2|k_1|\cdot t^{4/3}}}.
\end{equation}
A $C_c^\infty$ test function $f(t,\vec x) = \chi_{[\delta,2\delta]}(t)\,h(\vec x)$
with $\hat h(0)\neq 0$ then produces a $k_1$-IR divergence in the
smeared Wightman two-point:
\begin{align}
  \langle\phi(f)^2\rangle_\omega
  &\supset
  \int_{-K}^K dk_1\,|\hat h(k_1,0,0)|^2\,
   \biggl|\int_\delta^{2\delta}T_{\vec k}(t)\,dt\biggr|^2
  \\
  &\sim |\hat h(0)|^2 \int_{-K}^K \frac{dk_1}{|k_1|}\cdot O(\delta^{2/3})
  \\
  &= +\infty\quad(\text{any }\delta>0,\text{ at }K\text{ fixed}).
\end{align}
The IR divergence in $k_1$ is logarithmic and independent of the time-cutoff
$\delta$. (Sympy verification: \texttt{T2\_bianchi\_extension.py} block S3.)

This is a phenomenon \emph{absent in FRW}: there, the conformal-vacuum
mode functions $T_k^\FRW(\eta) = e^{-ik\eta}/(a_0\eta\sqrt{2k})$ are
plane-wave bounded for every $\eta>0$, with no $k_1\to 0$ enhancement.
The contracting Kasner direction is the algebraic source of the new
divergence, and it is intrinsically tied to the Weyl-non-vanishing
character of Bianchi I.

\paragraph{(iii) BFV folium uniformity.}

Either of (i) or (ii) alone gives an IR divergence of $\langle\phi(f)^2\rangle_\omega$
as $\delta\to 0^+$ for the SLE state $\omega$. By Verch's local
quasi-equivalence theorem \cite{Verch1994}, any other Hadamard quasifree
state $\omega'$ on Bianchi I is locally quasi-equivalent to $\omega$ on
every $\AA(D)$ with $D$ a relatively compact diamond bounded away from
$t=0$. Equivalence on the inductive limit follows from the
Brunetti--Fredenhagen--Verch generally covariant local quasi-equivalence
\cite{BFV2003}: the difference of two-point functions is a smooth integral
kernel, hence cannot cancel the $1/\sqrt{t_x t_y}$ singular prefactor or
the $1/|k_1|$ tachyonic IR divergence. Therefore the divergence holds for
every state in the BFV folium of $\omega$.

\paragraph{(iv) Conclusion.}

A cyclic-separating vector for $\AA(D_\BB)_\BI$ must, by construction,
carry a normal state on the inductive-limit algebra. Such a state would
have a finite smeared two-point on every test function in
$C_c^\infty(D_\BB)$, including those with support extending to $t=0$.
This contradicts (i), (ii), (iii). \qed

\section{Bianchi V and Bianchi IX coverage}

\begin{corollary}[T2 for Bianchi V, conditional]\label{cor:T2-BV}
Suppose a Hadamard quasifree state exists on the conformally coupled
massless scalar on vacuum Bianchi V (this is currently OPEN in the
AQFT-on-curved-spacetime literature; the most recent partial result is
Brum--Them \cite{Them-Brum2013} for inhomogeneous, compact Cauchy slices,
which does NOT directly cover non-compact $H^3$ Bianchi V). Then the
analog of Theorem~\ref{thm:T2-BI} holds, with the same proof structure:
the volume element $\sqrt{g_3}\sim t$ on the asymptotic Kasner attractor
\cite{WainwrightEllis1997}, and the contracting-direction tachyonic
divergence (ii) holds verbatim because the asymptotic Kasner exponents
of Bianchi V approach the same $(-1/3,2/3,2/3)$ point on the BKL
attractor.
\end{corollary}

\begin{question}[T2 for Bianchi IX / Mixmaster]\label{q:T2-BIX}
For Bianchi IX (Mixmaster) the past attractor is the chaotic Kasner-epoch
sequence (BKL 1970, \cite{BKL1970}), with the volume monotonically
vanishing $a_1 a_2 a_3 \to 0$ (Misner $\alpha\to-\infty$
\cite{Misner1969}) but the individual scale factors $a_i$ undergoing
unbounded chaotic bouncing. The volume-element argument (i) of
Theorem~\ref{thm:T2-BI} extends time-averagedly: the time-averaged
Hubble rate $\overline{H_{\text{avg}}} = \overline{\dot V/(3V)}$
is bounded below by a positive constant on the Kasner-attractor
invariant measure (Liouville measure on $\mathrm{PSL}(2,\Z)$ fundamental
domain; see Hartnoll--Yang \cite{HartnollYang2025} for the
modular-invariant quantum analog). Promoting this time-averaged bound
to a sample-pathwise smeared-two-point divergence requires a rigorous
invariant-measure analysis on the BKL attractor (Heinzle--Uggla 2009
established existence of a unique invariant measure on the attractor,
but the AQFT consequences have not been worked out). The
contracting-direction tachyonic divergence (ii) is more delicate:
during each Kasner epoch the contracting direction switches via the BKL
bounce maps, so there is no globally fixed contracting direction. A
Borel--Cantelli-type argument should give the divergence with probability
1 with respect to the BKL invariant measure, but a rigorous proof is
open.

\textbf{Conjecture.} Theorem~\ref{thm:T2-BI} extends to vacuum Bianchi IX
in the time-averaged sense, with probability 1 with respect to the
Heinzle--Uggla invariant measure on the BKL attractor. The Hartnoll--Yang
\cite{HartnollYang2025} ``end-of-time'' modular-invariant
Wheeler--DeWitt picture is in a structurally orthogonal framework
(WDW wavefunctional rather than local algebra) and cannot be plugged in
directly, but supplies a useful structural parallel: the modular-invariant
state at the BKL singularity is, in the WDW framework, an automorphic
$L$-function vanishing at the nontrivial Riemann zeros, which is the
WDW analog of the local-algebra non-cyclic-separating obstruction.
\end{question}

\section{What strategies S2, S4, S5 do (and do not) contribute}

We tested five strategies suggested by the parent agent. Theorem~\ref{thm:T2-BI}
is the result of S1$+$S3 (proof parts (i) and (ii) above). The remaining
three strategies are summarised here for completeness.

\subsection*{S2: DHN cosmological-billiards / Kac--Moody substitute}

Damour--Henneaux--Nicolai \cite{DHN2002} show that the asymptotic
near-singular dynamics of Einstein--matter in $D\ge 4$ is a billiard on
hyperbolic space, controlled by a hyperbolic Kac--Moody algebra ($E_{10}$
for 11d SUGRA, $AE_3$ for pure 4d gravity). This was hoped to provide a
``Kac--Moody substitute'' for the BFV folium structure used in the FRW T2
proof. We find this hope is structurally misplaced:
\begin{enumerate}
  \item The DHN dictionary is for the Wheeler--DeWitt wavefunctional on
  minisuperspace, not for the local algebra $\AA(D)_\BB$ on a fixed
  background. These are different mathematical objects.
  \item Hyperbolic Kac--Moody integrable highest-weight modules carry a
  positive-definite contravariant form (Kac \cite{Kac1990}, Ch.\ 11) and
  generate Type-I$_\infty$ von Neumann algebras (the universal
  enveloping algebra weakly-closed in a Type-I representation), NOT the
  Type-III/II structure required by our framework.
\end{enumerate}
\textbf{Verdict (S2):} structurally informative (the DHN/Hartnoll--Yang
near-singular structure is real and rich) but not a substitute for the
BFV folium argument.

\subsection*{S4: Tod isotropic-singularity contrapositive}

Tod \cite{Tod1987} characterised \emph{isotropic} singularities via a
Penrose conformal rescaling that regularises the Weyl tensor at $t=0$.
By contrapositive, vacuum Kasner (Weyl-divergent: $C_{txtx}\sim 1/t^{8/3}$,
sympy-verified) is NOT Tod-isotropic. This forbids the FRW T2
\emph{technique} (conformal pullback to Minkowski) but does not directly
forbid the FRW T2 \emph{conclusion} on Kasner; the conclusion requires
the volume-element / tachyonic-mode arguments above.

\textbf{Verdict (S4):} the contrapositive provides a clean DOMAIN-WALL
between Tod-isotropic singularities (handled by FRW T2 via conformal
pullback) and Weyl-divergent singularities (Bianchi I/V/IX, handled by
S1$+$S3 above). It is a CONSISTENCY CHECK for the framework, not a new
theorem.

\subsection*{S5: Apparent horizon entropy bound}

The cosmological apparent horizon for vacuum Kasner has comoving radius
$r_\mathrm{AH}=1/H_\mathrm{avg}=3t$, area $A_\mathrm{AH}=36\pi t^2$, and
Bekenstein--Bousso entropy bound $S_\mathrm{AH}=A_\mathrm{AH}/(4G_N) =
9\pi t^2/G_N \to 0$ as $t\to 0$. A vector that is cyclic-separating for
a non-trivial Type-III factor must have non-zero modular energy; if the
modular energy is bounded above by $S_\mathrm{AH}/(2\pi R)\to 0$, this
is a contradiction.

\textbf{Verdict (S5):} HEURISTIC, not rigorous: the semi-classical
Bousso bound assumes smooth-spacetime backreaction validity, which fails
exactly at the regime where the bound's conclusion ($A_\mathrm{AH}\to 0$)
is invoked. It supplies physical intuition for the result but cannot be
used as a proof step.

\section{Comparison with Penrose's original WCH}

Penrose's original Weyl Curvature Hypothesis \cite{Penrose1979,Penrose2010}
postulates that the Weyl tensor vanishes (or is bounded) at the past Big
Bang. In FRW (where $C_{abcd}\equiv 0$ by spatial isotropy), Penrose's
geometric content is automatically satisfied; the FRW T2 of
\cite{Remondiere2026FRW} is therefore a refinement of an empty
geometric statement into a non-trivial AQFT statement.

In Bianchi I/V/IX, Penrose's geometric WCH is non-trivial and FALSE: the
Weyl tensor diverges at the singularity (sympy-verified
\cite{RemondiereOpus2026Bianchi} CHECK 1). Theorem~\ref{thm:T2-BI} is
therefore a NEW algebraic statement that is independent of Penrose's
geometric criterion: it says that even when Penrose's WCH is violated
(divergent Weyl), the local algebra on the past-Big-Bang diamond
inductive limit fails to admit a cyclic-separating vector. The
algebraic asymmetry past$\leftrightarrow$future of the local-algebra
net is therefore robust to anisotropy, and is not contingent on
Penrose's WCH being true.

This is, to our knowledge, the first algebraic statement compatible
with the empirical fact that observed cosmological perturbations
$\delta T/T\sim 10^{-5}$ at last scattering are perfectly consistent with
either small-Weyl Penrose-WCH or with non-trivial-Weyl Bianchi-like
near-singular geometry: either way, the past-Big-Bang algebra is
non-cyclic-separating.

\section{Remaining gaps and outlook}

\begin{itemize}
  \item \textbf{Hadamard-state existence on Bianchi V/IX.} The most
  pressing technical gap. Banerjee--Niedermaier \cite{BanerjeeNiedermaier2023}
  cover Bianchi I; a corresponding construction on H$^3$ (Bianchi V) or
  on $S^3/\Gamma$ with Mixmaster dynamics (Bianchi IX) would allow
  Corollary~\ref{cor:T2-BV} and most of Question~\ref{q:T2-BIX} to be
  promoted to theorems. We estimate 6--12 months for an expert in
  microlocal analysis to extend the SLE construction to Bianchi V; the
  Bianchi IX case is a multi-year open problem entangled with the
  Heinzle--Uggla invariant-measure analysis.
  \item \textbf{Bianchi IX rigour.} The chaotic Kasner-epoch dynamics
  requires invariant-measure analysis on the BKL attractor to upgrade
  the time-averaged volume-element bound to a sample-pathwise
  divergence. Heinzle--Uggla 2009 established existence of an invariant
  measure; the AQFT consequences are open.
  \item \textbf{Connection to Hartnoll--Yang Wheeler--DeWitt picture.}
  The DHN/HY automorphic-$L$-function structure
  \cite{DHN2002,HartnollYang2025} is a structural parallel in a
  different framework. A precise dictionary between the WDW wavefunctional
  and our local-algebra net would clarify whether the modular invariance
  of the WDW state is the BB-singular limit of the BFV folium structure
  used here. This is a 2--5-year program.
\end{itemize}

\begin{thebibliography}{99}
\bibitem{Remondiere2026FRW}
K.\ Remondi\`ere, \emph{Algebraic Weyl Curvature Hypothesis: FRW Theorem T1+T2+T3},
companion note \texttt{algebraic\_arrow.tex} (May 2026).
\bibitem{RemondiereOpus2026Bianchi}
K.\ Remondi\`ere (with max-effort attempt by Opus 4.7),
\emph{Bianchi extension attempt: Proposition B.1 (T1 perturbative lift)},
companion note \texttt{bianchi\_extension\_opus.tex} (May 2026).
\bibitem{BanerjeeNiedermaier2023}
R.\ Banerjee and M.\ Niedermaier, \emph{States of low energy on Bianchi I spacetimes},
J.\ Math.\ Phys.\ \textbf{64}, 113503 (2023), arXiv:2305.11388.
\bibitem{Them-Brum2013}
K.\ Them and M.\ Brum, \emph{States of Low Energy in Homogeneous and Inhomogeneous,
Expanding Spacetimes}, Class.\ Quant.\ Grav.\ \textbf{30} (2013) 235035, arXiv:1302.3174.
\bibitem{Verch1994}
R.\ Verch, \emph{Local definiteness, primarity and quasi-equivalence of quasifree
Hadamard quantum states in curved spacetime}, Comm.\ Math.\ Phys.\ \textbf{160}
(1994) 507.
\bibitem{BFV2003}
R.\ Brunetti, K.\ Fredenhagen and R.\ Verch, \emph{The generally covariant locality
principle: A new paradigm for local quantum physics}, Comm.\ Math.\ Phys.\ \textbf{237}
(2003) 31, arXiv:math-ph/0112041.
\bibitem{HollandsWald2001}
S.\ Hollands and R.M.\ Wald, \emph{Local Wick polynomials and time-ordered products
of quantum fields in curved spacetime}, Comm.\ Math.\ Phys.\ \textbf{223} (2001) 289,
arXiv:gr-qc/0103074.
\bibitem{DHN2002}
T.\ Damour, M.\ Henneaux and H.\ Nicolai, \emph{Cosmological billiards}, Class.\
Quant.\ Grav.\ \textbf{20} (2003) R145, arXiv:hep-th/0212256.
\bibitem{DamourNicolai2004}
T.\ Damour and H.\ Nicolai, \emph{Eleven dimensional supergravity and the
$E_{10}/KE_{10}$ sigma-model at low $A_9$ levels}, arXiv:hep-th/0410245.
\bibitem{HartnollYang2025}
S.A.\ Hartnoll and M.\ Yang, \emph{The conformal primon gas at the end of time},
arXiv:2502.02661 (2025).
\bibitem{Tod1987}
K.P.\ Tod, \emph{Isotropic singularities and the $\gamma=2$ equation of state},
Class.\ Quant.\ Grav.\ \textbf{4} (1987) 1457; series of papers in the same
journal 1987--1992.
\bibitem{Kac1990}
V.\ Kac, \emph{Infinite Dimensional Lie Algebras}, 3rd ed., Cambridge University
Press 1990.
\bibitem{BKL1970}
V.A.\ Belinski, E.M.\ Lifshitz and I.M.\ Khalatnikov, \emph{Oscillatory
approach to a singular point in the relativistic cosmology}, Adv.\ Phys.\
\textbf{19} (1970) 525.
\bibitem{Misner1969}
C.W.\ Misner, \emph{Mixmaster universe}, Phys.\ Rev.\ Lett.\ \textbf{22}
(1969) 1071.
\bibitem{WainwrightEllis1997}
J.\ Wainwright and G.F.R.\ Ellis (eds.), \emph{Dynamical Systems in Cosmology},
Cambridge University Press 1997.
\bibitem{Penrose1979}
R.\ Penrose, \emph{Singularities and time-asymmetry}, in:
S.W.\ Hawking and W.\ Israel (eds.), \emph{General Relativity: An Einstein
Centenary Survey}, Cambridge University Press 1979.
\bibitem{Penrose2010}
R.\ Penrose, \emph{Cycles of Time}, Bodley Head 2010 (and arXiv:1011.3706).
\end{thebibliography}

\end{document}
